% AXIOM TEMPLATE: acm
% !! AXIOM-LOCKED-PREAMBLE-START !!  —— preamble 由模板锁定，agent 不得修改。
% 禁止改 \documentclass、禁止删除已有 \usepackage/\usetikzlibrary。
% Tectonic 会自动从 CTAN 下载 acmart.cls (首次编译联网)。
% format 选项可改: manuscript (默认,双栏) / acmsmall / acmlarge / acmtog / sigconf / sigplan
\documentclass[manuscript,review,anonymous]{acmart}

% —— acmart 自带: 数学 (amsmath/amssymb/bm)、graphicx、hyperref、natbib、cleveref 等 ——
% 注意: 不要重复 \usepackage{amssymb}/\usepackage{bm}, 会触发 \Bbbk already defined 冲突。
% —— 额外宏包 (acmart 未自带的) ——
\usepackage{mathtools}
\usepackage{booktabs}
\usepackage{multirow}
\usepackage{tabularx}

% —— TikZ (常用库全加载) ——
\usepackage{tikz}
\usetikzlibrary{positioning,arrows.meta,shapes.geometric,shapes.misc,calc,fit,backgrounds,decorations.pathreplacing}

% —— 伪代码 ——
\usepackage{algorithm}
\usepackage{algpseudocode}

% —— 列表 ——
\usepackage{enumitem}

% —— ACM 版权/会议信息 (投稿时按目标会议填) ——
\setcopyright{acmcopyright}
\acmConference{Conference Name}{Year}{Location}
\acmDOI{}
\acmISBN{}
% !! AXIOM-LOCKED-PREAMBLE-END !!

% AXIOM_BODY: agent 只能从 \begin{document} 之后开始编辑正文。
\begin{document}

\title{Paper Title}

\author{Author Name}
\email{email@example.com}
\affiliation{%
  \institution{Affiliation}
  \city{City}
  \country{Country}
}

\begin{abstract}
Replace this with the abstract.
\end{abstract}

\maketitle

\section{Introduction}
Replace this with the introduction.

% 引用: 写好 references.bib 后, 正文用 \cite{key}, 并取消下面两行注释。
% 未引用任何文献时保持注释, 否则空 .bbl 会触发 "missing \item" 报错。
% \bibliographystyle{ACM-Reference-Format}
% \bibliography{references}
\end{document}
